% ============================================================
%  Asfixia Neonatal / Encefalopatía Hipóxico-Isquémica
%  Artículo de revisión — formato working paper
% ============================================================
\documentclass[12pt]{article}

% ====== Página y formato básico ======
\usepackage[left=1.0in,right=1.0in,top=1.0in,bottom=1.0in]{geometry}
\usepackage{setspace}
\doublespacing
\usepackage{fancyhdr}
\pagestyle{fancy}
\fancyhf{}
\fancyfoot[C]{\thepage}
\renewcommand{\headrulewidth}{0pt}

% ====== Tipografía ======
\usepackage{lmodern}
\usepackage{microtype}
\usepackage[T1]{fontenc}

% ====== Matemáticas ======
\usepackage{amssymb, amsmath}

% ====== Tablas ======
\usepackage{array, booktabs, makecell}
\usepackage[flushleft]{threeparttable}

% ====== Figuras ======
\usepackage{graphicx}
\usepackage{caption}
\captionsetup{font=small, labelfont=bf, justification=justified}

% ====== Listas ======
\usepackage{enumitem}

% ====== URLs y colores ======
\usepackage{xurl}
\usepackage{xcolor}
\definecolor{citationcolor}{RGB}{0, 127, 255}

% ====== Bibliografía (biblatex + biber) ======
\usepackage[backend=biber,
            style=authoryear,
            maxcitenames=3,
            mincitenames=1,
            maxbibnames=99,
            giveninits=true,
            uniquename=false,
            uniquelist=true,
            dashed=false,
            natbib=true]{biblatex}
\addbibresource{asfixia_neonatal.bib}
\renewcommand*{\nameyeardelim}{\addcomma\space}

% ====== Idioma ======
\usepackage[spanish]{babel}

% ====== Hyperref (último) ======
\usepackage[hidelinks, breaklinks, colorlinks=true,
            linkcolor=citationcolor, citecolor=citationcolor,
            urlcolor=citationcolor]{hyperref}

% ============================================================
%  Metadatos
% ============================================================
\title{Asfixia Neonatal y Encefalopatía Hipóxico-Isquémica:\\
       Fisiopatología, Diagnóstico, Tratamiento y Pronóstico%
       \thanks{Los autores no declaran conflictos de interés.
               Este trabajo no recibió financiamiento externo.}}

\author{
  Author One\thanks{Departamento de Neonatología, Facultad de Medicina,
  Universidad Ficticia. Correo: author.one@univficticia.edu} \quad
  Author Two\thanks{Unidad de Cuidados Intensivos Neonatales,
  Hospital Ficticio Nacional. Correo: author.two@hospitalficticio.gob}
}

\date{Abril de 2026}

% ============================================================
\begin{document}
% ============================================================

\maketitle
\thispagestyle{empty}

% ============================================================
\begin{abstract}
\noindent \singlespacing
\textbf{Introducción.} La asfixia neonatal representa la segunda causa de años de vida
ajustados por discapacidad en el período neonatal y genera una carga desproporcionada
en países de ingresos bajos y medianos.
\textbf{Métodos.} Revisión narrativa de la literatura clínica y epidemiológica publicada
entre 1976 y 2024, con énfasis en ensayos clínicos aleatorizados, metaanálisis y guías
de práctica clínica iberoamericanas.
\textbf{Resultados.} La hipotermia terapéutica (33,5\,°C rectal durante 72 horas),
iniciada antes de las 6 horas de vida en recién nacidos con encefalopatía
hipóxico-isquémica moderada o grave, reduce la mortalidad y la discapacidad
neuológica mayor con un número necesario a tratar de 6 a 9.
El diagnóstico se apoya en criterios clínicos estandarizados, electroencefalografía
de amplitud integrada y resonancia magnética.
\textbf{Conclusión.} La detección oportuna y el inicio precoz de la hipotermia
constituyen la intervención más eficaz disponible; su acceso universal, especialmente
en entornos de recursos limitados, sigue siendo el principal desafío.
\end{abstract}

\vspace{1em}
\noindent \textbf{Códigos JEL:} I10, I12, I18

\vspace{0.5em}
\noindent \textbf{Palabras clave:} asfixia neonatal, encefalopatía hipóxico-isquémica,
hipotermia terapéutica, fisiopatología neonatal, pronóstico neurológico

\newpage
\setcounter{page}{1}

% ============================================================
\section{Introducción}
\label{sec:introduccion}
% ============================================================

La asfixia neonatal ocurre entre 1 y 6 por cada 1\,000 nacidos vivos en países de
ingresos altos, cifra que asciende a 25 por cada 1\,000 en contextos de bajos recursos;
en Ecuador, datos del Ministerio de Salud Pública registraron 7 casos por cada 100
nacidos vivos en 2022 \parencite{ecuador2022}. El 96\,\% de la carga global de esta
condición recae sobre países de ingresos bajos y medianos \parencite{kurinczuk2010},
lo que convierte a la asfixia neonatal en un problema de equidad sanitaria tanto como
en un desafío clínico.

La magnitud de sus consecuencias justifica la atención. La asfixia neonatal causa el
23\,\% de la mortalidad infantil global y explica el 20\,\% de los casos de parálisis
cerebral infantil \parencite{lawn2005}. En términos de carga de enfermedad, constituye
la segunda causa de años de vida ajustados por discapacidad en el período neonatal,
superada únicamente por la prematuridad \parencite{kurinczuk2010}. La mortalidad no es
el único desenlace relevante: los supervivientes pueden presentar epilepsia, deterioro
cognitivo y trastornos motores que persisten a lo largo de toda la vida.

La encefalopatía hipóxico-isquémica (EHI) es la expresión clínica más grave de la
asfixia neonatal. Se define como el conjunto de signos neurológicos que aparecen en las
primeras horas de vida como consecuencia de la reducción aguda del aporte de oxígeno
y substrato energético al cerebro \parencite{volpe2012}. La distinción entre asfixia
neonatal ---evento bioquímico y clínico de privación de oxígeno--- y EHI ---su repercusión
cerebral--- tiene implicancias prácticas: no todos los recién nacidos asfixiados
desarrollan EHI, pero quienes lo hacen concentran el mayor riesgo de muerte y
discapacidad.

El único tratamiento con eficacia demostrada en ensayos clínicos aleatorizados es la
hipotermia terapéutica \parencite{shankaran2005, gluckman2005, azzopardi2008}. Su
ventana de acción es estrecha ---las primeras seis horas de vida---, lo que hace del
reconocimiento precoz una condición sine qua non para su aplicación efectiva. El
presente artículo revisa la fisiopatología, los criterios diagnósticos, el protocolo de
tratamiento y los factores pronósticos de la EHI, con atención a las particularidades
relevantes para entornos latinoamericanos.

% ============================================================
\section{Fisiopatología}
\label{sec:fisiopatologia}
% ============================================================

El daño cerebral en la EHI no ocurre en un único momento sino en una cascada que se
desarrolla en cuatro fases, cada una con características bioquímicas y temporales
propias \parencite{ferriero2004}. Comprender esta secuencia es fundamental para
entender por qué la hipotermia funciona y cuál es su ventana terapéutica.

\textbf{Fase I: evento centinela.} La reducción aguda del flujo sanguíneo
---por desprendimiento de placenta, prolapso de cordón, paro cardiorrespiratorio u
otro evento--- suprime la entrega de oxígeno y glucosa. Las células cerebrales pasan
del metabolismo aeróbico al anaeróbico, con producción de lactato y caída del ATP.
La falla energética activa los canales iónicos voltaje-dependientes y permite la
entrada masiva de calcio al citoplasma, desencadenando la excitotoxicidad primaria.
Este proceso se inicia en minutos y puede terminar en necrosis celular si la
isquemia es prolongada.

\textbf{Fase II: período de latencia.} La restauración de la circulación ---espontánea
o mediante reanimación--- inaugura un intervalo de relativa estabilidad metabólica que
dura entre 1 y 6 horas. La actividad eléctrocerebral se recupera parcialmente,
el metabolismo oxidativo mejora y el daño visible es todavía mínimo. Este período
de latencia constituye la ventana terapéutica: es el momento en que la hipotermia
puede interrumpir la cascada \parencite{edwards2010}. Pasadas las 6 horas, la
eficacia del tratamiento disminuye de manera sustancial.

\textbf{Fase III: fase secundaria.} A partir de las 6 horas, y con un pico entre las
24 y las 72 horas, se activan cuatro mecanismos que amplifican el daño neuronal.
La excitotoxicidad secundaria es mediada por glutamato y potenciada por la
acumulación de calcio intracelular. El estrés oxidativo genera especies reactivas de
oxígeno que atacan lípidos, proteínas y ADN. La respuesta inflamatoria recluta
microglia y macrófagos circulantes que liberan citocinas proinflamatorias
(IL-1$\beta$, IL-6, TNF-$\alpha$). Por último, la apoptosis ---activada por las vías
intrínseca y extrínseca de la caspasa--- elimina neuronas que podrían haberse
recuperado si el proceso hubiese sido interrumpido \parencite{volpe2012}. La
hipotermia actúa principalmente suprimiendo estos cuatro mecanismos secundarios.

\textbf{Fase IV: fase terciaria.} En las semanas y meses posteriores al evento
agudo persisten fenómenos de remodelación, neuroinflamación crónica y epigenética
alterada que pueden continuar generando daño \parencite{ferriero2004}. Esta fase
justifica la investigación de neuroprotectores adyuvantes ---eritropoyetina, melatonina,
células madre--- cuyo beneficio en combinación con la hipotermia se evalúa actualmente
en ensayos clínicos.

La asfixia no afecta solo al cerebro. La redistribución del gasto cardíaco ante la
hipoxia ---respuesta de buceo--- privilegia el flujo coronario y cerebral a expensas de
órganos esplácnicos, riñón y piel. El resultado es una disfunción multiorgánica cuya
severidad es paralela a la del daño cerebral: cardiopatía transitoria con hipertensión
pulmonar, insuficiencia renal aguda, enterocolitis necrosante, elevación de
transaminasas e hipoglucemia son hallazgos frecuentes que requieren monitoreo y
tratamiento activo \parencite{shankaran2005}. La presencia y magnitud de la afectación
sistémica sirven como indicador barrométrico de la gravedad global del cuadro.

% ============================================================
\section{Diagnóstico}
\label{sec:diagnostico}
% ============================================================

El diagnóstico de EHI requiere la confluencia de criterios clínicos, bioquímicos y
neurofisiológicos \parencite{alsepneo2021}. La clasificación en grados de severidad
determina directamente la elegibilidad para hipotermia terapéutica, por lo que una
evaluación rigurosa y sistematizada es condición necesaria para el tratamiento.

\subsection{Criterios diagnósticos}
\label{subsec:criterios}

Los criterios diagnósticos se organizan en dos grupos. Los \textbf{criterios A}
documentan la ocurrencia y la severidad del evento hipóxico-isquémico: presencia de
evento centinela identificable, más al menos uno de los siguientes ---puntuación de
Apgar $\leq 5$ a los cinco minutos de vida, necesidad de reanimación cardiopulmonar
avanzada por más de 10 minutos, pH en sangre de cordón $< 7{,}0$, o déficit de
bases $\geq 16$\,mmol/L \parencite{sarnat1976}. Los \textbf{criterios B} documentan
la afectación neurológica mediante escalas estandarizadas de encefalopatía: la escala
de Sarnat modificada clasifica la EHI en grados I (leve), II (moderada) y III (grave)
con base en el nivel de conciencia, tono muscular, reflejos primitivos y presencia de
convulsiones \parencite{sarnat1976}. La escala cuantitativa de García-Alix asigna una
puntuación que estratifica como leve ($< 8$ puntos), moderada (8--30 puntos) y grave
($> 30$ puntos), con mayor objetividad interobservador \parencite{garciaalix2012}.
La elegibilidad para hipotermia exige que ambos grupos de criterios estén presentes y
que la EHI sea moderada o grave.

\subsection{Monitoreo neurofisiológico}
\label{subsec:eeg}

El electroencefalograma de amplitud integrada (aEEG) permite la monitorización
continua del patrón electrocortical con un equipo de cabecera y escasa formación
técnica. Sus patrones ---fondo continuo, discontinuo, burst-suppression, voltaje bajo
continuo o isoeléctrico--- tienen valor diagnóstico para estratificar la severidad y
pronóstico para predecir el desenlace neurológico \parencite{alsepneo2021}. Su
principal utilidad práctica es la detección de convulsiones subclínicas, que ocurren
hasta en el 50\,\% de los neonatos con EHI moderada-grave y son clínicamente silentes
\parencite{volpe2012}. El EEG convencional de 19 canales aporta mayor sensibilidad
diagnóstica y debe utilizarse ante cualquier duda sobre la interpretación del aEEG o
cuando se sospeche estado epiléptico.

\subsection{Neuroimagen}
\label{subsec:rmn}

La resonancia magnética (RMN) es la herramienta pronóstica principal en la EHI.
Su rendimiento diagnóstico varía con el momento de adquisición. La RMN precoz
($< 5$ días), con secuencias de difusión (DWI) y coeficiente de difusión aparente
(ADC), detecta restricción de difusión en las regiones lesionadas ---particularmente
ganglios basales, tálamo y corteza--- antes de que las secuencias convencionales
muestren anormalidades \parencite{ferriero2004}. La RMN tardía (5--14 días), con
secuencias T1 y T2, refleja con mayor fidelidad la extensión final del daño y
constituye el mejor predictor imagenológico del pronóstico motor y cognitivo. La
lesión bilateral de ganglios basales y tálamo en T1 predice parálisis cerebral
discinética con una probabilidad que varía del 10 al 98\,\% según la extensión
\parencite{garciaalix2012}.

\subsection{Biomarcadores}
\label{subsec:biomark}

La enolasa neuronal específica (NSE) es el biomarcador sérico más estudiado en la EHI.
Concentraciones superiores a 50\,ng/ml durante las primeras 24 horas se asocian con
mal pronóstico neurológico con alta especificidad, aunque su sensibilidad es moderada
\parencite{kurinczuk2010}. Otros marcadores ---S100B, proteína de unión a ácidos grasos
de tipo cerebral (FABP), neurofilamento de cadena ligera--- se encuentran en evaluación
clínica pero carecen aún de umbrales de decisión validados para uso rutinario en
entornos iberoamericanos.

% ============================================================
\section{Tratamiento}
\label{sec:tratamiento}
% ============================================================

\subsection{Hipotermia terapéutica}
\label{subsec:hipotermia}

La hipotermia terapéutica es el estándar de oro para el tratamiento de la EHI moderada
y grave. Tres ensayos aleatorizados multicéntricos ---CoolCap \parencite{gluckman2005},
NICHD \parencite{shankaran2005} y TOBY \parencite{azzopardi2008}--- demostraron
reducción significativa de la mortalidad y la discapacidad mayor, con un número
necesario a tratar (NNT) de 6 a 9. El metaanálisis de \citet{edwards2010} confirmó
que el enfriamiento a 33--34\,°C durante 72 horas reduce en un 25\,\% el riesgo
combinado de muerte o discapacidad neuológica grave a los 18 meses.

El protocolo estándar establece como temperatura objetivo 33,5\,°C rectal, mantenida
durante 72 horas sin interrupción \parencite{jacobs2013}. El recalentamiento posterior
es lento y controlado: 0,3--0,5\,°C por hora para evitar el rebote convulsivo e
hipotérmico. El enfriamiento puede ser pasivo ---retirar las fuentes de calor y dejar
que el recién nacido se enfríe de manera natural--- o activo mediante mantas o equipos
de enfriamiento servo-controlados. En entornos de recursos limitados, el enfriamiento
pasivo o con bolsas de gel frío es una alternativa que las guías latinoamericanas
aceptan como medida puente \parencite{alsepneo2021}.

Las indicaciones precisas son: edad gestacional $\geq 35$ semanas, peso al nacer
$\geq 1{,}800$\,g, edad postnatal $\leq 6$ horas, y EHI moderada o grave documentada
con criterios A y B \parencite{statpearls2024}. Las contraindicaciones absolutas
incluyen malformaciones congénitas mayores incompatibles con la vida, coagulopatía
grave no corregible y situaciones en que el tratamiento no esté en sintonía con los
objetivos de cuidado acordados con la familia. La prematuridad extrema
($< 35$ semanas) es contraindicación relativa por la falta de datos de seguridad
en este subgrupo.

El tiempo de inicio es crítico: datos del ensayo NICHD y del estudio TOBY indican que
el inicio de la hipotermia antes de los 180 minutos de vida se asocia con mejores
resultados motores que su inicio entre los 180 y 360 minutos \parencite{shankaran2014}.
Este gradiente temporal refuerza la necesidad de protocolos de identificación precoz
en sala de partos y de traslados neonatales organizados para garantizar que el recién
nacido llegue al centro de hipotermia dentro de la ventana.

\subsection{Monitoreo durante la hipotermia}
\label{subsec:monitoreo}

El manejo durante las 72 horas de hipotermia requiere vigilancia sistémica y
neurológica continua. Los parámetros mínimos son: glucemia cada 2--4 horas (objetivo
$> 47$\,mg/dL), electrolitos séricos, coagulograma, función hepática y renal al inicio
y cada 24 horas, aEEG o EEG continuo para detección de convulsiones, y ecocardiografía
para descartar hipertensión pulmonar persistente y disfunción miocárdica
\parencite{alsepneo2021}. La hipotermia altera el metabolismo de múltiples fármacos;
las dosis de sedación, antiepilépticos y antibióticos deben ajustarse.

\subsection{Manejo de convulsiones}
\label{subsec:convulsiones}

Las convulsiones neonatales complican hasta el 40--45\,\% de los casos de EHI y
agravan el daño cerebral al incrementar la demanda metabólica en el contexto de
reserva energética reducida \parencite{volpe2012}. El fenobarbital en dosis de carga
de 20\,mg/kg intravenoso es el fármaco de primera línea; puede repetirse hasta
alcanzar 40\,mg/kg si las convulsiones persisten \parencite{alsepneo2021}. El
levetiracetam ---dosis de 20--40\,mg/kg intravenoso--- se reserva como tercera línea por
su perfil de seguridad y la menor carga de efectos adversos sedantes, aunque la
evidencia comparativa directa con otros fármacos en neonatos es todavía limitada
\parencite{statpearls2024}. La fenitoína y el midazolam se utilizan como segunda línea
en los esquemas de la mayoría de los centros latinoamericanos.

% ============================================================
\section{Pronóstico y seguimiento}
\label{sec:pronostico}
% ============================================================

El pronóstico de la EHI depende de la severidad del episodio, del patrón de lesión en
la RMN y de la respuesta al tratamiento, y abarca un espectro amplio que va desde la
recuperación neurológica completa hasta el fallecimiento en período neonatal.

\subsection{Secuelas neurológicas}
\label{subsec:secuelas}

La lesión bilateral de ganglios basales y tálamo ---patrón más frecuente en la EHI
grave--- se asocia con parálisis cerebral (PC) discinética. La probabilidad de PC varía
del 10 al 98\,\% según la extensión de la lesión en la RMN \parencite{garciaalix2012},
lo que ilustra la importancia de la neuroimagen como herramienta pronóstica individual
y no solo poblacional. La epilepsia se presenta en el 40--45\,\% de los casos de EHI
con convulsiones neonatales y puede requerir tratamiento antiepiléptico crónico
\parencite{volpe2012}. Entre los supervivientes con PC, las comorbilidades son la norma:
dolor crónico en el 75\,\%, discapacidad intelectual en el 50\,\% y trastornos
visuales ---incluyendo cortical visual impairment--- en el 50\,\% de los casos
\parencite{kurinczuk2010}.

\subsection{Factores pronósticos}
\label{subsec:factores}

Cuatro factores tienen el mayor peso pronóstico individual. Primero, la severidad
clínica al ingreso medida con la escala de Sarnat: los neonatos con EHI grado III
tienen mortalidad hospitalaria superior al 50\,\% y, entre los supervivientes, más del
70\,\% presenta discapacidad mayor \parencite{sarnat1976}. Segundo, el patrón de la
RMN: la lesión de ganglios basales/tálamo confiere peor pronóstico que la lesión
predominantemente cortical \parencite{ferriero2004}. Tercero, el tiempo al inicio de
la hipotermia: cada hora de retraso dentro de la ventana de 6 horas se asocia con
peor outcome \parencite{shankaran2014}. Cuarto, la persistencia de convulsiones a
pesar del tratamiento antiepiléptico es un marcador de gravedad independiente
\parencite{volpe2012}.

\subsection{Seguimiento a largo plazo}
\label{subsec:seguimiento}

Las guías de la ALSEPNEO \parencite{alsepneo2021} recomiendan seguimiento
neuropediátrico mínimo hasta los seis años de edad. Este seguimiento debe incluir
evaluación del neurodesarrollo con instrumentos estandarizados ---Bayley Scales of
Infant and Toddler Development, escala de Griffiths--- cada seis meses hasta los dos
años y anualmente hasta los seis. La detección precoz de déficits motores, cognitivos
o sensoriales permite iniciar intervenciones terapéuticas ---fisioterapia, terapia
ocupacional, estimulación temprana--- en el período de mayor plasticidad neuronal.
El seguimiento continuado es también necesario para documentar de manera prospectiva
los resultados de los programas de hipotermia, lo que permite actualizar los datos
de eficacia en contextos locales donde la incidencia de comorbilidades prenatales
difiere de los entornos de los ensayos originales.

La tabla~\ref{tab:pronostico} sintetiza los principales desenlaces según la
clasificación de Sarnat al ingreso.

\begin{table}[ht]
\centering
\begin{threeparttable}
\caption{Pronóstico según grado de encefalopatía hipóxico-isquémica al ingreso}
\label{tab:pronostico}
\begin{tabular}{lccc}
\toprule
\textbf{Desenlace} & \textbf{EHI grado I} & \textbf{EHI grado II} & \textbf{EHI grado III} \\
\midrule
Mortalidad neonatal & $< 5$\,\% & 15--25\,\% & $> 50$\,\% \\
Discapacidad mayor & $< 5$\,\% & 25--35\,\% & $> 70$\,\% \\
Parálisis cerebral & $< 1$\,\% & 10--20\,\% & 40--60\,\% \\
Epilepsia & $< 5$\,\% & 20--30\,\% & 30--50\,\% \\
\bottomrule
\end{tabular}
\begin{tablenotes}
\small
\item \textit{Nota.} Estimaciones basadas en \citet{sarnat1976} y \citet{kurinczuk2010},
con ajustes según la revisión de \citet{edwards2010} que incluye datos post-hipotermia.
EHI = encefalopatía hipóxico-isquémica. Los valores expresan rangos aproximados;
la variabilidad individual es considerable y debe completarse con la evaluación
por RMN. Muestra: ensayos TOBY, NICHD y CoolCap combinados ($n > 600$).
\end{tablenotes}
\end{threeparttable}
\end{table}

% ============================================================
\section{Conclusiones}
\label{sec:conclusiones}
% ============================================================

La asfixia neonatal y su expresión clínica principal, la EHI, configuran un problema
de salud pública de magnitud relevante tanto en términos de mortalidad como de carga
de discapacidad, con una distribución geográfica que concentra el 96\,\% del daño en
países de ingresos bajos y medianos. En Ecuador y en el conjunto de América Latina,
la incidencia supera la de los países de donde provienen la mayoría de los ensayos
clínicos, lo que refuerza la urgencia de adaptar los protocolos diagnósticos y
terapéuticos a la realidad local.

El conocimiento acumulado sobre la fisiopatología en cuatro fases ofrece una base
racional para la hipotermia terapéutica, cuya eficacia está respaldada por evidencia
de nivel I procedente de tres ensayos aleatorizados independientes y ratificada en el
metaanálisis de Edwards y colaboradores, con un NNT de 6 a 9 para prevenir muerte o
discapacidad mayor. El tiempo al inicio del tratamiento es el modificador más potente
de su eficacia: cada hora dentro de la ventana de 6 horas cuenta, y los sistemas de
salud que no disponen de unidades de hipotermia deben organizar circuitos de traslado
neonatal que garanticen el acceso oportuno.

Los avances diagnósticos ---aEEG, RMN y biomarcadores--- permiten hoy estratificar el
riesgo con mayor precisión que hace una década, aunque persiste la necesidad de
validar umbrales y protocolos en poblaciones latinoamericanas. La investigación en
neuroprotectores adyuvantes ---eritropoyetina, melatonina, sulfato de magnesio--- abre
perspectivas para mejorar los resultados más allá de lo alcanzable con la hipotermia
sola, pero requiere ensayos adicionales antes de su adopción rutinaria.

El seguimiento neuropediátrico hasta los seis años es tan parte del tratamiento como
la hipotermia misma: sin él, las secuelas se detectan tarde, las intervenciones de
habilitación se retrasan y se pierde la oportunidad de actuar durante el período
de mayor plasticidad cerebral. Garantizar ese seguimiento en sistemas de salud con
alta dispersión geográfica y escasos especialistas representa el reto inmediato para
los equipos de neonatología en la región.

% ============================================================
\clearpage
\small
\singlespacing
\printbibliography[title={Referencias}]
% ============================================================

\end{document}
